\section{Introduction}
\label{sec:intro}

When a production issue is reported, an engineer typically receives two
things: a natural-language description of what went wrong, and logs. There
is no failing test, no debugger session, and often no way to reproduce the
problem. From this evidence alone, the engineer must find the code that
caused it. The challenge this paper answers calls for techniques that
perform this step automatically: given an issue report, logs, and a source
code snapshot, produce the affected module, a ranked list of suspicious
code regions, and the evidence that connects them.

The core difficulty is that logs show where a failure was \emph{detected},
not where it was \emph{caused}. A real example from our benchmark makes
this concrete. Kafka issue KAFKA-6624 reports that a disk is wrongly
marked offline, and its stack trace points at
\texttt{FileChannelImpl.force} and \texttt{FileRecords.flush}, the code
that hit a closed file channel. The actual fix landed in
\texttt{LogManager.scala}, which allowed a log to be deleted while it was
still being flushed. Every file the trace names is innocent. Any method
that works by matching the report against the code it mentions, which is
what an engineer grepping the codebase does, starts from the wrong files.
In our benchmark, \nHardPct{} of real reports name none of the files the
fix later touched.

Existing techniques approach this from two sides, and both stop short.
Classical fault localization ranks suspicious code using failing tests or
coverage, execution evidence that does not exist at report
time~\cite{just2014defects4j}. Recent agent-based variants retrieve over
the codebase itself~\cite{rafi2024multiagent}, treating every issue as if
it were the project's first. Log-based root cause analysis in cloud
operations~\cite{wang2024rcagent} identifies the faulty service or
component, but stops above the level of code regions.

One source of evidence is missing from both: the project's own resolved
history. Software failures recur. When an issue arrives, some past issue
has often failed in a similar way, and its fix commit records exactly
which files were at fault that time. We call this \emph{incident
memory}. Past incidents whose reports resemble the new one vote for the
files their fixes touched. A memory used carelessly would cheat, because
the history of a project contains fixes that did not exist when the
report was filed, so every use of it here sits behind a temporal
firewall that admits only incidents resolved before that moment.

Studying any of this required a benchmark that was not available: issue
reports with logs, the code as it was when the issue was open, and
ground-truth faulty files. We mined one from two public datasets and
project git histories. It contains \nInstances{} real incidents from four
projects of four organizations (Apache Kafka, MariaDB, MongoDB, Qt),
where the ground truth of each instance is the set of source files its
fix commit modified, and the snapshot is the parent of that commit. The
benchmark is stratified by exactly the difficulty the challenge names:
whether the report names any faulty file.

We then use the benchmark to compare four families of method on the same
instances, under one protocol and one firewall: classic
information-retrieval bug localization, lexical log-to-code retrieval
fused with incident memory, a general-purpose open 27B model prompted
with the report and the file tree, and that model used to re-rank
retrieved candidates. Accuracy in the top five rises across the families
from $0.180$ to $0.839$. The more useful result is that no family
dominates. Retrieval alone and the model alone solve different
instances: retrieval is the only one of the two to succeed on 6.1\% of
instances, the model alone on 21.1\%, and the union of their top-five
lists reaches $0.892$ where each alone reaches $0.681$ and $0.832$.

That gap between what the parts find separately and what either finds
alone is the opening this paper takes. We present \sys{}, an agent that
keeps the retrieval and the incident memory as tools and adds tools that
read code, resolve a symbol from a stack trace to the file that defines
it, and consult past incidents behind the same firewall. \sys{} runs on
the same 27B model used in the study, so what it adds is the
architecture, not a larger model. Kafka issue KAFKA-7693 shows the
pattern. The report is a crash on a negative sequence number, and
retrieval returns nothing useful. \sys{} resolves the symbol in the
trace, reads the code that produces that number, and ranks
\texttt{TransactionManager.java} first, the file whose arithmetic
creates the negative value. On 105 of the 375 hard instances it ranks the
faulty file first where retrieval does not place it in the top five at
all.

This paper makes four contributions.
\begin{itemize}
\item \textbf{A benchmark} of \nInstances{} issue--log--code--fix
  instances across four organizations, with a per-filter datasheet and a
  per-project incident-memory pool, released for reuse
  (\cref{sec:benchmark}).
\item \textbf{A comparative study} of four families of localization
  method on those instances, under one protocol and one temporal
  firewall (\cref{sec:results}). It reports what each family achieves,
  where each fails, and that they fail on different instances. It also
  measures what the firewall is worth: the classic similar-report
  signal, evaluated without report-time admissibility, overstates Acc@5
  by 9\% to 43\% relative (\cref{sec:results:leakage}).
\item \textbf{\sys{}}, an agent that turns the retrieval, the incident
  memory, and the code itself into tools (\cref{sec:agent}). It reaches
  Acc@5 $0.903$ and Acc@1 $0.723$, improving on every family in the
  study and on the best combination of them, with the widest margins on
  the reports that name no faulty file. Its cost is reported next to its
  accuracy.
\item \textbf{Three validity checks} that a reviewer of an
  LLM-based result should expect, each of which could have invalidated
  the result and did not (\cref{sec:results:validity}). The model cannot
  recall these public fixes from the issue key alone. Shuffling the
  order of the retrieved candidates changes no outcome. Every baseline
  is rescored on the exact instances the agent ran, because the
  evaluation sample is easier than the full split.
\end{itemize}

The lexical parts of the system are deterministic, and the agent's tools
are pure functions of the snapshot, so a recorded run replays to
identical tool results. Every reported number traces to a result file
through a released script.
